\textbf{Approved title pair (synthetic)}\par\smallskip
{\footnotesize\setlength{\tabcolsep}{4pt}\renewcommand{\arraystretch}{1.0}
\begin{tabular}{@{}>{\raggedright\arraybackslash}p{0.24\linewidth}>{\raggedright\arraybackslash}p{0.72\linewidth}@{}}
\texttt{label\_source} & human\_approved\_title\_pair \\
\texttt{domain} & labour \\
\texttt{content\_overlap} & 0.286 \\
\texttt{question} & \bng{অন্যায় আচরণের জন্য শাস্তি সংক্রান্ত একটা সমস্যায় পড়েছি, কোথায় যাব}? \\
\texttt{positive} & \bng{বাংলাদেশ ইপিজেড শ্রম আইন, ২০১৯}, \S \bng{১৫১} \bng{অন্যায় আচরণের জন্য দণ্ড} \\
\texttt{hard negative 1} & The Prisoners Act, 1900, \S 30 Lunatic prisoners how to be dealt with. \\
\texttt{hard negative 2} & The Railways Act, 1890, \S 113 Travelling without pass or ticket or with insufficient pass or ticket of beyond authorised distance \\
\texttt{hard negative 3} & \bng{নির্যাতন এবং হেফাজতে মৃত্যু (নিবারণ) আইন, ২০১৩}, \S \bng{১৬} \bng{আপীল} \\
\end{tabular}}
\bigskip
\textbf{Authored narrative (synthetic)}\par\smallskip
{\footnotesize\setlength{\tabcolsep}{4pt}\renewcommand{\arraystretch}{1.0}
\begin{tabular}{@{}>{\raggedright\arraybackslash}p{0.24\linewidth}>{\raggedright\arraybackslash}p{0.72\linewidth}@{}}
\texttt{label\_source} & authored\_against\_provision\_text \\
\texttt{domain} & family \\
\texttt{content\_overlap} & 0.136 \\
\texttt{question} & \bng{স্বামীর বিরুদ্ধে খোরপোশের মামলা করেছি}। \bng{গত তারিখে আমি খুব অসুস্থ থাকায় যেতে পারিনি, আমার উকিলও সেদিন অন্য আদালতে আটকে গিয়েছিলেন}। \bng{পরে শুনলাম স্বামীও সেদিন যায়নি}। \bng{এখন মামলার কী হবে}? \\
\texttt{positive} & \bng{পারিবারিক আদালত আইন, ২০২৩}, \S \bng{১০} \bng{পক্ষগণের অনুপস্থিতির ফলাফল} \\
\texttt{hard negative 1} & \bng{আইন-শৃঙ্খলা বিঘ্নকারী অপরাধ (দ্রুত বিচার) আইন, ২০০২}, \S \bng{১১} \bng{আসামীর অনুপস্থিতিতে বিচার} \\
\texttt{hard negative 2} & \bng{মানব পাচার ও অভিবাসী চোরাচালান প্রতিরোধ ও দমন আইন, ২০২৬}, \S \bng{৪৩} \bng{আসামির অনুপস্থিতিতে বিচার} \\
\texttt{hard negative 3} & The Code of Criminal Procedure, 1898, \S 199B  \\
\end{tabular}}
